\documentclass[11pt]{article}
\usepackage[margin=1in]{geometry}
\usepackage{amsmath,amssymb,amsthm}
\usepackage{hyperref}

\newtheorem{theorem}{Theorem}
\newtheorem{lemma}{Lemma}
\newtheorem{proposition}{Proposition}

\title{A Counterexample to General Convergence of Double Q-Learning with Linear Function Approximation}
\author{}
\date{}

\begin{document}
\maketitle

\section{Claim}

Double Q-learning is convergent in the tabular case under standard assumptions. This proof shows that the same is not true for arbitrary linear function approximation.

\begin{theorem}
There exists a finite discounted Markov decision process, a stationary behavior policy, a bounded linear feature map, and a Robbins--Monro stepsize sequence such that symmetric Double Q-learning with linear function approximation is unbounded almost surely from any strictly positive initialization. This holds even though the optimal action-value function is identically zero and is exactly representable by the feature map.
\end{theorem}

The example has one state, two actions, zero rewards, and one scalar feature.

\section{The MDP, features, and behavior policy}

Consider the discounted MDP with one state $s$, two actions $a_1, a_2$, deterministic self-transition, zero rewards, and discount factor $\gamma = 0.9$. Thus $P(s \mid s,a) = 1$ and $r(s,a) = 0$ for both actions $a \in \{a_1, a_2\}$. Since all rewards are zero,
\begin{equation}\label{eq:qstar}
Q^*(s,a_1) = Q^*(s,a_2) = 0.
\end{equation}

Use one-dimensional linear features $\phi(s,a_1) = 1$ and $\phi(s,a_2) = 2$. The two Double-Q estimators are parameterized by scalars $\theta_A, \theta_B$:
\begin{equation}\label{eq:qa-qb}
Q_A(s,a) = \theta_A \phi(s,a), \qquad Q_B(s,a) = \theta_B \phi(s,a).
\end{equation}
The optimal function $Q^* \equiv 0$ is exactly representable by the parameter value $\theta_A = \theta_B = 0$.

Let the stationary behavior policy be $\mu(a_1 \mid s) = 0.9$ and $\mu(a_2 \mid s) = 0.1$. At time $n$, sample an action from $\mu$, and write $X_{n+1} = \phi(s, A_{n+1}) \in \{1, 2\}$. Hence
\begin{equation}\label{eq:feature-dist}
\mathbb{P}(X_{n+1} = 1) = 0.9, \qquad \mathbb{P}(X_{n+1} = 2) = 0.1.
\end{equation}

Independently of the sampled action and of the past, choose estimator $A$ or estimator $B$ for update with probability $1/2$ each. Use the global stepsize sequence
\begin{equation}\label{eq:stepsize}
\alpha_n = \frac{a}{n+1}, \qquad 0 < a \leq \frac{1}{16}.
\end{equation}
Then $\sum_{n=0}^{\infty} \alpha_n = \infty$ and $\sum_{n=0}^{\infty} \alpha_n^2 < \infty$. Initialize with arbitrary strictly positive parameters: $\theta_{A,0} > 0$ and $\theta_{B,0} > 0$.

\section{The Double Q-learning updates}

We use the standard symmetric Double Q-learning update. If estimator $A$ is selected at time $n$, then
\begin{equation}\label{eq:update-a}
\theta_{A,n+1}
= \theta_{A,n}
+ \alpha_n \phi(s, A_{n+1})
\bigl[\gamma\, Q_B\bigl(s, \arg\max_{a'} Q_A(s,a')\bigr) - Q_A(s, A_{n+1})\bigr],
\end{equation}
while $\theta_B$ is unchanged. If estimator $B$ is selected, the roles are reversed:
\begin{equation}\label{eq:update-b}
\theta_{B,n+1}
= \theta_{B,n}
+ \alpha_n \phi(s, A_{n+1})
\bigl[\gamma\, Q_A\bigl(s, \arg\max_{a'} Q_B(s,a')\bigr) - Q_B(s, A_{n+1})\bigr],
\end{equation}
while $\theta_A$ is unchanged. Because the state never changes and the reward is zero, these are the full updates.

\section{Positivity is invariant}

We first show that the positive quadrant is invariant. This will imply that the greedy action is fixed throughout the entire trajectory.

Suppose $\theta_A > 0$. Then $Q_A(s,a_2) = 2\theta_A > \theta_A = Q_A(s,a_1)$, so the greedy action under $Q_A$ is uniquely $a_2$. Similarly, if $\theta_B > 0$, the greedy action under $Q_B$ is uniquely $a_2$. Thus no tie-breaking convention is needed as long as the parameters remain strictly positive.

Assume $\theta_A, \theta_B > 0$, and let $X \in \{1, 2\}$ be the sampled feature. If estimator $A$ is updated, then the target action is $a_2$, whose feature is $2$. Therefore
\begin{equation}\label{eq:update-expanded}
\theta_A^+ = \theta_A + \alpha X(\gamma \cdot 2\theta_B - X\theta_A)
= (1 - \alpha X^2)\theta_A + 1.8\,\alpha X\,\theta_B.
\end{equation}
Because $X \leq 2$ and $\alpha_n \leq 1/16$, we have $1 - \alpha_n X^2 \geq 1 - 4\alpha_n \geq 3/4 > 0$. Thus $\theta_A^+ > 0$. The same argument applies when estimator $B$ is updated:
\begin{equation}
\theta_B^+ = (1 - \alpha X^2)\theta_B + 1.8\,\alpha X\,\theta_A > 0.
\end{equation}
Therefore, since $\theta_{A,0} > 0$ and $\theta_{B,0} > 0$, we have $\theta_{A,n} > 0$ and $\theta_{B,n} > 0$ for every $n$. Consequently, the greedy action is always $a_2$ for both estimators. Along this trajectory, the nonlinear maximization in the Double-Q target reduces to a fixed linear target.

\section{Dynamics of the parameter sum}

Define $S_n = \theta_{A,n} + \theta_{B,n}$. By positivity invariance, $S_n > 0$ for all $n$. We will prove that $S_n \to \infty$ almost surely. This implies that the Double-Q parameter vector is unbounded and therefore cannot converge to $(0,0)$, which represents $Q^*$.

Let $r_n = \theta_{A,n} / S_n \in (0,1)$. If estimator $A$ is updated at time $n$, then
\begin{equation}
S_{n+1} = S_n + \alpha_n X_{n+1}(1.8\,\theta_{B,n} - X_{n+1}\,\theta_{A,n}).
\end{equation}
Dividing by $S_n$, this becomes $S_{n+1} = S_n(1 + \alpha_n Y_{n+1})$, where, in the case of an $A$-update,
\begin{equation}\label{eq:y-a-update}
Y_{n+1} = X_{n+1}\bigl(1.8(1 - r_n) - X_{n+1}\,r_n\bigr).
\end{equation}
If estimator $B$ is updated, then similarly $S_{n+1} = S_n(1 + \alpha_n Y_{n+1})$, where
\begin{equation}\label{eq:y-b-update}
Y_{n+1} = X_{n+1}\bigl(1.8\,r_n - X_{n+1}(1 - r_n)\bigr).
\end{equation}
For both possible updates, because $r_n \in (0,1)$ and $X_{n+1} \in \{1, 2\}$, we have $-4 \leq Y_{n+1} \leq 3.6$. Therefore $|\alpha_n Y_{n+1}| \leq 1/4$ and $1 + \alpha_n Y_{n+1} \geq 3/4 > 0$. Thus the logarithm of $S_n$ is well defined for every $n$.

\section{Positive drift of the logarithm}

Let $\mathcal{F}_n$ be the history up to time $n$, including $\theta_{A,n}$ and $\theta_{B,n}$, but not the fresh action sample or estimator choice at time $n$. Conditional on $\mathcal{F}_n$, the ratio $r_n$ is fixed.

First average over the independent choice of which estimator is updated. For a fixed sampled feature $X$,
\begin{align}
\mathbb{E}[Y_{n+1} \mid \mathcal{F}_n, X_{n+1} = X]
&= \tfrac{1}{2}\, X\bigl(1.8(1 - r_n) - X\,r_n\bigr)
 + \tfrac{1}{2}\, X\bigl(1.8\,r_n - X(1 - r_n)\bigr) \notag \\
&= \tfrac{1}{2}\, X(1.8 - X). \label{eq:cond-mean-x}
\end{align}
Now average over the behavior policy:
\begin{align}
\mathbb{E}[Y_{n+1} \mid \mathcal{F}_n]
&= 0.9 \cdot \tfrac{1}{2} \cdot 1 \cdot (1.8 - 1)
 + 0.1 \cdot \tfrac{1}{2} \cdot 2 \cdot (1.8 - 2) \notag \\
&= 0.9 \cdot 0.4 + 0.1 \cdot (-0.2) \notag \\
&= 0.36 - 0.02 = 0.34. \label{eq:drift}
\end{align}
Thus $Y_{n+1}$ has a strictly positive conditional mean, uniformly over the current parameters.

Define $L_n = \log S_n$. Since $S_{n+1} = S_n(1 + \alpha_n Y_{n+1})$, we have $L_{n+1} - L_n = \log(1 + \alpha_n Y_{n+1})$. For $|t| \leq 1/4$, the elementary inequality $\log(1 + t) \geq t - t^2$ holds. Applying this with $t = \alpha_n Y_{n+1}$, and using $|Y_{n+1}| \leq 4$, gives
\begin{equation}\label{eq:log-drift-bound}
\mathbb{E}[L_{n+1} - L_n \mid \mathcal{F}_n]
\geq \alpha_n\,\mathbb{E}[Y_{n+1} \mid \mathcal{F}_n] - \alpha_n^2\,\mathbb{E}[Y_{n+1}^2 \mid \mathcal{F}_n]
\geq 0.34\,\alpha_n - 16\,\alpha_n^2.
\end{equation}
Because $\sum_{n=0}^{\infty} \alpha_n = \infty$ and $\sum_{n=0}^{\infty} \alpha_n^2 < \infty$, we have $\sum_{n=0}^{N-1}(0.34\,\alpha_n - 16\,\alpha_n^2) \to +\infty$. So the predictable drift of $L_n$ has a deterministic lower bound whose partial sums diverge to $+\infty$.

\section{Martingale fluctuations are negligible}

Let $\Delta L_{n+1} = L_{n+1} - L_n = \log(1 + \alpha_n Y_{n+1})$, and define the martingale difference
\begin{equation}
M_{n+1} = \Delta L_{n+1} - \mathbb{E}[\Delta L_{n+1} \mid \mathcal{F}_n].
\end{equation}
For $|t| \leq 1/4$, there is a universal constant $K$ such that $|\log(1+t)| \leq K|t|$. For example, one may take $K = 2$. Therefore $|\Delta L_{n+1}| \leq 2\alpha_n |Y_{n+1}| \leq 8\alpha_n$. Consequently,
\begin{equation}
\bigl|\mathbb{E}[\Delta L_{n+1} \mid \mathcal{F}_n]\bigr|
\leq \mathbb{E}\bigl[|\Delta L_{n+1}| \mid \mathcal{F}_n\bigr]
\leq 8\alpha_n.
\end{equation}
Hence $|M_{n+1}| \leq |\Delta L_{n+1}| + |\mathbb{E}[\Delta L_{n+1} \mid \mathcal{F}_n]| \leq 16\alpha_n$. It follows that $\mathbb{E}[M_{n+1}^2 \mid \mathcal{F}_n] \leq 256\,\alpha_n^2$. Since $\sum_n \alpha_n^2 < \infty$,
\begin{equation}
\sum_{n=0}^{\infty} \mathbb{E}[M_{n+1}^2 \mid \mathcal{F}_n] < \infty.
\end{equation}
By the martingale convergence theorem for square-integrable martingales with summable conditional variances, the martingale partial sums $\sum_{n=0}^{N-1} M_{n+1}$ converge almost surely to a finite random limit.

\section{Almost-sure divergence}

Using the decomposition
\begin{equation}\label{eq:doob}
L_N = L_0 + \sum_{n=0}^{N-1} \mathbb{E}[\Delta L_{n+1} \mid \mathcal{F}_n] + \sum_{n=0}^{N-1} M_{n+1},
\end{equation}
and the lower bound from Section~6,
\begin{equation}
\sum_{n=0}^{N-1} \mathbb{E}[\Delta L_{n+1} \mid \mathcal{F}_n]
\geq \sum_{n=0}^{N-1} (0.34\,\alpha_n - 16\,\alpha_n^2) \to +\infty.
\end{equation}
The martingale term converges almost surely to a finite limit. Therefore $L_N \to +\infty$ almost surely. Since $L_N = \log S_N$, this implies $S_N \to \infty$ almost surely. Since both parameters remain positive,
\begin{equation}
\max\{\theta_{A,N},\, \theta_{B,N}\} \geq \frac{S_N}{2} \to \infty.
\end{equation}
Thus the parameter vector is unbounded almost surely. In particular, the Double-Q iterates cannot converge to the representable optimum $(\theta_A, \theta_B) = (0,0)$. This proves the theorem. \qed

\section{Relation to the expected linear dynamics}

This section is not needed for the almost-sure proof, but it explains the source of the instability. In the positive quadrant, the greedy action is always $a_2$. Therefore the expected update is linear. Let $C = \mathbb{E}[X^2]$ and $D = \gamma\,\mathbb{E}[2X]$. In this example, $C = 0.9 \cdot 1^2 + 0.1 \cdot 2^2 = 1.3$ and $D = 0.9(0.9 \cdot 2 \cdot 1 + 0.1 \cdot 2 \cdot 2) = 1.98$. The expected Double-Q update has the form
\begin{equation}
\mathbb{E}[\Delta \theta_A] = \frac{\alpha}{2}(D\,\theta_B - C\,\theta_A),
\qquad
\mathbb{E}[\Delta \theta_B] = \frac{\alpha}{2}(D\,\theta_A - C\,\theta_B).
\end{equation}
Equivalently,
\begin{equation}\label{eq:matrix-dynamics}
\mathbb{E}
\begin{bmatrix}
\Delta \theta_A \\
\Delta \theta_B
\end{bmatrix}
= \frac{\alpha}{2}
\begin{bmatrix}
-C & D \\
D & -C
\end{bmatrix}
\begin{bmatrix}
\theta_A \\
\theta_B
\end{bmatrix}.
\end{equation}
Introduce the average and difference coordinates $u = (\theta_A + \theta_B)/2$ and $v = (\theta_A - \theta_B)/2$. The corresponding mean ODE decouples as
\begin{equation}\label{eq:ode}
\dot{u} = \tfrac{1}{2}(D - C)\,u, \qquad \dot{v} = -\tfrac{1}{2}(C + D)\,v.
\end{equation}
Since $D - C = 1.98 - 1.3 = 0.68 > 0$, we obtain $\dot{u} = 0.34\,u$. Thus the average mode is unstable. The almost-sure argument above turns this mean-dynamics instability into a genuine stochastic nonconvergence result.

\section{What the counterexample shows and does not show}

This counterexample shows that there cannot be a general convergence theorem for Double Q-learning with arbitrary linear function approximation under off-policy sampling and bootstrapping.

It does not show that Double Q-learning always diverges. Double Q-learning is convergent in tabular settings under the usual assumptions, and it may converge in restricted linear settings with additional structure. The point is that Double Q-learning does not, by itself, remove the instability caused by the combination of bootstrapping, off-policy sampling, and function approximation.

The example is deliberately minimal: one state, two actions, zero rewards, one scalar feature, and a fixed behavior policy. Despite this simplicity, the Double-Q parameter norm diverges almost surely from every strictly positive initialization.

\end{document}
